\begin{figure}[htbp]
\centering
\begin{tikzpicture}[x=1cm,y=1cm,font=\small,
  active/.style={line width=2pt,draw=blue!65!black},
  event/.style={densely dotted,draw=black!24},
  endpoint/.style={circle,draw=blue!65!black,fill=white,
    line width=.65pt,inner sep=0pt,minimum size=4pt},
  endpoint label/.style={font=\footnotesize,fill=white,inner sep=2pt},
  energy/.style={line width=1pt,draw=black},
  jump/.style={-{Stealth[length=3pt]},densely dashed,
    line width=.8pt,draw=red!65!black}]
  \foreach \t in {2.3,3.8,6.4,9.0,11.6}
    \draw[event] (\t,.1)--(\t,6.15);

  \node[anchor=west] at (1.9,6.65) {Output-dependent scale intervals};
  \foreach \yy in {5.8,4.8,3.8}
    \draw[black!22] (1.95,\yy)--(12.25,\yy);
  \node[anchor=east] at (1.7,5.8) {$(i_1,j_1)$};
  \node[anchor=east] at (1.7,4.8) {$(i_2,j_2)$};
  \node[anchor=east] at (1.7,3.8) {$(i_3,j_3)$};
  \draw[active] (2.3,5.8)--(6.4,5.8);
  \draw[active] (3.8,4.8)--(9.0,4.8);
  \draw[active] (6.4,3.8)--(11.6,3.8);
  \foreach \xx/\yy in {2.3/5.8,6.4/5.8,3.8/4.8,9.0/4.8,6.4/3.8,11.6/3.8}
    \node[endpoint] at (\xx,\yy) {};
  \node[endpoint label,above] at (2.3,5.91) {$\varepsilon_{i_1j_1}^{\,2}$};
  \node[endpoint label,above] at (6.4,5.91) {$\rho_{i_1j_1}^{\,2}$};
  \node[endpoint label,above] at (3.8,4.91) {$\varepsilon_{i_2j_2}^{\,2}$};
  \node[endpoint label,above] at (9.0,4.91) {$\rho_{i_2j_2}^{\,2}$};
  \node[endpoint label,above] at (6.4,3.91) {$\varepsilon_{i_3j_3}^{\,2}$};
  \node[endpoint label,above] at (11.6,3.91) {$\rho_{i_3j_3}^{\,2}$};
  \draw[-{Stealth[length=4pt]},black!65] (1.95,3.35)--(12.6,3.35)
    node[right] {$s$};

  \node[anchor=west,fill=white,inner sep=2pt] at (1.9,2.9) {Integrated energy};
  \draw[-{Stealth[length=4pt]},black!65] (1.95,.1)--(12.6,.1)
    node[right] {$s$};
  \node[anchor=east] at (1.7,1.7) {$J(s)$};
  \draw[energy] (2.3,2.2) .. controls (2.9,2.05) and (3.3,1.91) .. (3.8,1.82);
  \draw[jump] (3.8,1.82)--(3.8,2.38);
  \draw[energy] (3.8,2.38) .. controls (4.7,2.13) and (5.4,1.85) .. (6.4,1.64);
  \draw[jump] (6.4,1.64)--(6.4,1.97);
  \draw[energy] (6.4,1.97) .. controls (7.2,1.75) and (8.1,1.37) .. (9.0,1.18);
  \draw[jump] (9.0,1.18)--(9.0,.78);
  \draw[energy] (9.0,.78) .. controls (9.9,.61) and (10.7,.43) .. (11.6,.32);
  \node[below] at (2.3,.02) {$s_-$};
  \node[below] at (11.6,.02) {$s_+$};
\end{tikzpicture}
\caption{Three selected output pairs and a schematic integrated energy,
not an exact heat-flow solution. The entry $A_0^s(i,j)$ is retained only
on $(\varepsilon_{ij}^2,\rho_{ij}^2)$. Between successive event scales,
$J(s)$ is nonincreasing by the fixed-array dissipation inequality.
Dashed arrows connect its one-sided limits at events; the values at the
events themselves are not shown. Jumps may have either sign, and
endpoints may coincide, as illustrated
by $\rho_{i_1j_1}^2=\varepsilon_{i_3j_3}^2$. The bound
$\sum_\tau |\Delta J(\tau)|\le C\mathcal S(A)$, together with the
initial-energy bound, controls the accumulated dissipation independently
of the number of event scales.}
\label{fig:mask-energy}
\end{figure}
